\section{Discussion}
\label{sec:discussion}

The evidence presented in the preceding sections suggests that 
rebound effects represent a structural challenge to the 
assumption that AI-driven efficiency gains will reduce global 
CO\textsubscript{2} emissions. However, several qualifications 
are necessary before drawing broader conclusions.

First, the argument advanced in this paper should not be 
interpreted as a wholesale rejection of AI's potential 
contribution to climate change mitigation. A substantial body 
of research identifies high-impact areas in which machine 
learning can meaningfully reduce greenhouse gas emissions, 
including the optimization of electricity grids, improvements 
in climate modelling, and the development of alternative 
materials \cite{rolnick2022}. The question is not whether AI 
can deliver efficiency gains, but whether those gains are 
sufficient to offset the resource consumption they induce.

Second, applying the Jevons Paradox to AI involves conceptual 
limitations. The paradox was originally formulated in the 
context of a single commodity -- coal -- and its explanatory 
power diminishes when applied to the complex, multi-sectoral 
dynamics of AI diffusion \cite{luccioni2025}. The causal 
relationships between AI efficiency, economic expansion, and 
behavioral responses are difficult to isolate and may vary 
significantly across sectors and regions \cite{luccioni2025}.

Third, the empirical evidence presented here is subject to 
geographic constraints. The threshold analysis of the digital 
economy's effect on energy consumption is based on panel data 
from 30 Chinese provinces and may not be directly generalizable 
to other economic contexts \cite{chen2025}. Furthermore, the 
absence of standardized reporting from AI companies makes 
comprehensive lifecycle assessments difficult to conduct 
\cite{luccioni2025}.

These limitations notwithstanding, the structural pattern 
identified across multiple sources is consistent: efficiency 
improvements in AI systems tend to stimulate demand rather 
than reduce it. Addressing this dynamic requires more than 
technical optimization. Effective policy responses would 
need to combine efficiency measures with regulatory frameworks 
that constrain overall resource consumption, decouple AI 
development from unsustainable growth, and mandate transparent 
reporting of environmental impacts \cite{luccioni2025}. As 
long as AI development operates within market logics that 
reward scale and growth, efficiency gains are likely to 
function as enablers of expansion rather than substitutes 
for it \cite{widdicks2023}.